% ============ §13 Graded rings, the Hilbert function and the Samuel function (原书 pp.92-101) ============
\section{Graded rings, the Hilbert function and the Samuel function}\label{sec:13}\index{graded rings, section 13}\index{Hilbert function, section 13}\index{Samuel function, section 13}

Let $G$ be an Abelian semigroup with identity element $0$; (that is, $G$
is a set with an addition law $+$ satisfying associativity $(x + y) + z = x +
(y + z)$, commutativity $x + y = y + x$, and such that $0 + x = x$). A \emph{graded} (or
$G$-graded) ring is a ring $R$ together with a direct sum decomposition of $R$ as
an additive group $R = \bigoplus_{i \in G} R_i$ satisfying $R_iR_j \subset R_{i+j}$. Similarly, a \emph{graded}
$R$-module is an $R$-module $M$ together with a direct sum decomposition $M =
\bigoplus_{i \in G} M_i$ satisfying $R_iM_j \subset M_{i+j}$. An element $x \in M$ is \emph{homogeneous}\index{homogeneous element} if
$x \in M_i$ for some $i \in G$, and $i$ is then called the \emph{degree} of $x$. A general element
$x \in M$ can be written uniquely in the form $x = \sum_{i \in G} x_i$ with $x_i \in M_i$ and only
finitely many $x_i \neq 0$; $x_i$ is called the \emph{homogeneous term} of $x$ of degree $i$.

A submodule $N \subset M$ is called a \emph{homogeneous submodule}\index{homogeneous submodule} (or \emph{graded}
submodule) if it can be generated by homogeneous elements. This
condition is equivalent to either of the following two:
\begin{enumerate}[(1)]
  \item For $x \in M$, if $x \in N$ then each homogeneous term of $x$ is in $N$;
\end{enumerate}
For a homogeneous submodule $N \subset M$ we set $N_i = N \cap M_i$; then
$M/N = \bigoplus_{i \in G} M_i/N_i$ is again a graded $R$-module.

One sees from the definition that $R_0 \subset R$ is a subring, and that each
graded piece $M_i$ of a graded $R$-module $M$ is an $R_0$-module.

The notion of graded ring is most frequently used when $G$ is the
semigroup $\{0, 1, 2, \ldots\}$ of non-negative integers, which we denote by $\N$. In this
case, we set $R^+ = \bigoplus_{n > 0} R_n$; then $R^+$ is an ideal of $R$, with $R/R^+ \cong R_0$.

The polynomial ring $R = R_0[X_1, \ldots, X_r]$ over a ring $R_0$ is usually made
into an $\N$-graded ring by defining the degree of a monomial $X_1^{\alpha_1} \cdots X_r^{\alpha_r}$ as
the total degree $\alpha_1 + \cdots + \alpha_r$; however, $R$ has other useful gradings. For
example, $R$ has an $\N^r$-grading in which $X_1^{\alpha_1} \cdots X_r^{\alpha_r}$ has degree $(\alpha_1, \ldots, \alpha_r)$;
the value of systematically using this grading can be seen in Goto--
Watanabe [1]. Alternatively, giving each of the $X_i$ some suitable weight
$d_i$ and letting the monomial $X_1^{\alpha_1} \cdots X_r^{\alpha_r}$ have weight $\sum \alpha_i d_i$ defines an
$\N$-grading of $R$. For example, the ring $R_0[X, Y, Z]/(f)$, where $f = a_1X^{\alpha_1} +
a_2Y^{\alpha_2} + a_3Z^{\alpha_3}$ can be graded by giving the images of $X, Y, Z$ the weights
$\alpha_1, \alpha_2$ and $\alpha_3$, respectively.

A \emph{filtration}\index{filtration} of a ring $A$ is a descending chain $A = J_0 \supset J_1 \supset \cdots$ of ideals
such that $J_nJ_m \subset J_{n+m}$; the associated graded ring $\gr(A)$ is defined as
follows. First of all as a module we set $\gr_n(A) = J_n/J_{n+1}$ for $n \ge 0$, and
$\gr(A) = \bigoplus_{n \ge 0} \gr_n(A)$; then we define the product by
\[
  (x + J_{n+1})(y + J_{m+1}) = xy + J_{n+m+1} \quad \text{for } x \in J_n \text{ and } y \in J_m.
\]
It is easy to see that $\gr(A)$ becomes a graded ring. The filtration
$J_1 \supset J_2 \supset \cdots$ defines a linear topology on $A$ (see \S\ref{sec:8}), and the completion $\hat{A}$ of
$A$ in this topology has a filtration $\hat{J}_1 \supset \hat{J}_2 \supset \cdots$ such that $\hat{A}/\hat{J}_n \cong A/J_n$ for
all $n$, hence $\hat{J}_n/\hat{J}_{n+1} \cong J_n/J_{n+1}$, and
\[
  \gr(\hat{A}) = \gr(A).
\]

Let $A$ be a ring, $I$ an ideal, and let $B = \bigoplus_{n \ge 0} I^n/I^{n+1}$ be the graded ring
associated with the filtration $I \supset I^2 \supset \cdots$ of $A$ by powers of $I$; the various
notations $\gr_I(A)$, $\gr^I(A)$ and $G_I(I)$ are used to denote $B$ in the current
literature. An element of $B_n = I^n/I^{n+1}$ can be expressed as a linear
combination of products of $n$ elements of $B_1 = I/I^2$, so that $B$ is generated
over the subring $B_0 = A/I$ by elements of $B_1$. If $I = Ax_1 + \cdots + Ax_r$ and $\bar{x}_i$
denotes the image of $x_i$ in $B_1 = I/I^2$ then
\[
  B = \gr_I(A) = (A/I)[\bar{x}_1, \ldots, \bar{x}_r],
\]
and $B$ is a quotient of the polynomial ring $(A/I)[X_1, \ldots, X_r]$ as a graded
ring.

\begin{theorem}\label{thm:13.1}
An $\N$-graded ring $R = \bigoplus_{n \ge 0} R_n$ is Noetherian if and only if
$R_0$ is Noetherian and $R$ is finitely generated as a ring over $R_0$.
\end{theorem}

\begin{proof}
The `if' is obvious, and we prove the `only if': suppose that $R$
is Noetherian. Then since $R_0 \cong R/R^+$, $R_0$ is Noetherian. $R^+$ is a home-
geneous ideal, and is finitely generated, so that we can suppose that it is
generated by homogeneous elements $x_1, \ldots, x_r$. Then it is easy to see that
$R = R_0[x_1, \ldots, x_r]$; in fact it is enough to show that $R_n \subset R_0[x_1, \ldots, x_r]$
for every $n$. Now writing $d_i$ for the degree of $x_i$ we have
\[
  (*) \quad R_n = x_1R_{n-d_1} + x_2R_{n-d_2} + \cdots + x_rR_{n-d_r}.
\]
Indeed, for $y \in R_n$, write $y = \sum c_ix_i$ with $f_i \in R$; then setting $g_i$ for the
homogeneous term of degree $n - d_i$ of $f_i$ (with $g_i = 0$ if $n - d_i < 0$), we also
have $y = \sum x_ig_i$. From $(*)$ it follows by induction that $R_n \subset R_0[x_1, \ldots, x_r]$.
\end{proof}

Let $R = \bigoplus_{n \ge 0} R_n$ be a Noetherian graded ring; then if $M = \bigoplus_{n \ge 0} M_n$
is a finitely generated graded $R$-module, each $M_n$ is finitely generated
as $R_0$-module. In fact when $M = R$ this is clear from $(*)$ above. In the general
case $M$ can be generated by a finite number of homogeneous elements $\omega_i$: $M
= \sum R\omega_i$. Now letting $e_i$ be the degree of $\omega_i$, we have as above that
\[
  M_n = R_{n-e_1}\omega_1 + \cdots + R_{n-e_r}\omega_r \quad \text{(where } R_i = 0 \text{ for } i < 0),
\]
and hence $M_n$ is a finite $R_0$-module. In particular if $R_0$ is an Artinian ring,
then $l(M_n) < \infty$, where $l$ denotes the length of an $R_0$-module. In this case we
define the \emph{Hilbert series}\index{Hilbert series} $P(M, t)$ of $M$ by the formula:
\[
  P(M, t) = \sum_{n=0}^{\infty} l(M_n)t^n \in \Z[[t]].
\]
(In combinatorics it is a standard procedure to associate with a sequence of
numbers $a_0, a_1, a_2, \ldots$ the generating function $\sum a_it^i$.)

\begin{theorem}\label{thm:13.2}
Let $R = \bigoplus_{n \ge 0} R_n$ be a Noetherian graded ring with $R_0$
Artinian, and let $M$ be a finitely generated graded $R$-module. Suppose that
$R = R_0[x_1, \ldots, x_r]$ with $x_i$ of degree $d_i$, and that $P(M, t)$ is as above. Then
$P(M, t)$ is a rational function of $t$, and can be written
\[
  P(M, t) = f(t) \Big/ \prod_{i=1}^{r} (1 - t^{d_i}),
\]
where $f(t)$ is a polynomial with coefficients in $\Z$.
\end{theorem}

\begin{proof}
By induction on $r$, the number of generators of $R$. When $r = 0$ we
have $R = R_0$, so that for $n$ sufficiently large, $M_n = 0$, and the power series
$P(M, t)$ is a polynomial. When $r > 0$, multiplication by $x_r$ defines an $R_0$-
linear map $M_n \to M_{n+d_r}$; writing $K_n$ and $L_{n+d_r}$ for the kernel and cokernel,
we get an exact sequence
\[
  0 \to K_n \to M_n \to M_{n+d_r} \to L_{n+d_r} \to 0.
\]
Set $K = \bigoplus K_n$, and $L = \bigoplus L_n$. Then $K$ is a submodule of $M$, and $L = M/x_rM$,
so that $K$ and $L$ are finite $R$-modules; moreover $x_rK = x_rL = 0$, so that $K$
and $L$ can be viewed as $R/x_rR$-modules, and hence we can apply the
induction hypothesis to $P(K, t)$ and $P(L, t)$. Now from the above exact
sequence we get
\[
  l(K_n) - l(M_n) + l(M_{n+d_r}) - l(L_{n+d_r}) = 0.
\]
If we multiply this by $t^{n+d_r}$ and sum over $n$ this gives
\[
  t^{d_r}P(K, t) - t^{d_r}P(M, t) + P(M, t) - P(L, t) = g(t),
\]
where $g(t) \in \Z[t]$. The theorem follows at once from this.
\end{proof}

A lot of information on the values of $l(M_n)$ can be obtained from the
above theorem. Especially simple is the case $d_1 = \cdots = d_r = 1$, so that $R$ is
generated over $R_0$ by elements of degree $1$. In this case $P(M, t) = f(t)
(1 - t)^{-r}$; if $f(t)$ has $(1 - t)$ as a factor we can cancel to get $P$ in the form
\[
  P(M, t) = f(t)(1 - t)^{-d} \quad \text{with } f \in \Z[t], \; d \ge 0,
\]
and if $d > 0$ then $f(1) \neq 0$.
If this holds, we will write $d = d(M)$. Since $(1 - t)^{-1} = 1 + t + t^2 + \cdots$, we can
repeatedly differentiate both sides to get
\[
  (1 - t)^{-d} = \sum_{m \ge 0} \binom{m + d - 1}{d - 1} t^m.
\]
(This can of course be proved in other ways, for example by induction on $d$.)
If $f(t) = a_0 + a_1t + \cdots + a_st^s$ then
\[
  (*) \quad l(M_n) = a_0\binom{n+d-1}{d-1} + a_1\binom{n+d-2}{d-1} + \cdots + a_s\binom{n-s+d-1}{d-1},
\]
where we interpret $\binom{m}{d-1} = 0$ for $m < d - 1$. The right-hand side of $(*)$ can be
formally rearranged as a polynomial in $n$ with rational coefficients, say $\varphi(n)$;
then
\[
  \varphi(X) = \frac{f(1)}{(d-1)!}X^{d-1} + \text{(terms of lower degree)}.
\]
Since $\binom{m + d - 1}{d-1}$ coincides with the polynomial $m(m - 1) \cdots (m - d + 2)/
(d-1)!$ for $m \ge 0$, this implies the following result.

\begin{corollary*}
If $d_1 = \cdots = d_r = 1$ in Theorem 2, and $d = d(M)$ is defined
as above, then there is a polynomial $\varphi_M(X)$ of degree $d - 1$ with rational
coefficients such that for $n \ge s + 1 - d$ we have $l(M_n) = \varphi_M(n)$. Here $s$ is the
degree of the polynomial $(1 - t)^dP(M, t)$.
\end{corollary*}

The polynomial $\varphi_M$ appearing here is called the \emph{Hilbert polynomial}\index{Hilbert polynomial} of the
graded module $M$. The numerical function $l(M_n)$ itself is called the \emph{Hilbert}
\emph{function} of $M$; by the degree of a Hilbert function we mean the degree of the
corresponding Hilbert polynomial.

\begin{remark}
For general $d_1, \ldots, d_r$, it is no longer necessarily the case that $l(M_n)$
can be represented by one polynomial.
\end{remark}

\begin{example}\label{ex:13.1}
When $R = R_0[X_1, \ldots, X_r]$, the number of monomials of
degree $n$ is $\binom{n+r-1}{r-1}$, so that
\[
  l(R_n) = l(R_0)\binom{n+r-1}{r-1}
\]
holds for every $n \ge 0$, and the right-hand side is $\varphi_R(n)$. Thus $\varphi_R(X) =
(l(R_0)/r!)(X + r)(X + r - 1) \cdots (X + 1)$.
\end{example}

\begin{example}\label{ex:13.2}
Let $k$ be a field, and $F(X_1, \ldots, X_r)$ a homogeneous poly-
nomial of degree $s$; set $R = k[X_1, \ldots, X_r]/(F(X))$. Then for $n \ge s$,
\[
  l(R_n) = \binom{n+r-1}{r-1} - \binom{n-s+r-1}{r-1},
\]
and hence, setting $\binom{n+r-1}{r-1} = (1/r!)n^{r-1} + a_1n^{r-2} + \cdots$, we have
\[
  \varphi_R(X) = \frac{1}{r!}[X^{r-1} - (X - s)^{r-1}] + a_1[X^{r-2} - (X - s)^{r-2}] + \cdots
\]
\[
  = \frac{s}{(r-1)!}X^{r-2} + \text{(terms of lower degree)}.
\]
\end{example}

\begin{example}\label{ex:13.3}
Let $k$ be a field, and $R = k[X_1, \ldots, X_r]/P = k[\xi_1, \ldots, \xi_r]$,
where $P$ is a homogeneous prime ideal. Let $t$ be the transcendence degree
of $R$ over $k$, and suppose that $\xi_1, \ldots, \xi_t$ are algebraically independent
over $k$; then there are $\binom{n+t-1}{t-1}$ monomials of degree $n$ in the
$\xi_1, \ldots, \xi_t$, and these are linearly independent over $k$, so that
\[
  l(R_n) \ge \binom{n+t-1}{t-1},
\]
from which it follows that $d \ge t$. In fact we will prove
later (Theorem 8) that $d = t$.
\end{example}

A homogeneous ideal of the polynomial ring $k[X_1, \ldots, X_r]$ over a field
$k$ defines an algebraic variety in $r$-dimensional projective space $\mathbb{P}^r$, and
the Hilbert polynomial plays an important role in algebraic geometry.
For example, note that the numerator of the leading term of $\varphi_R$ in Example
2 is equal to the degree of $F$. This holds in more generality, but we must
leave details of this to textbooks on algebraic geometry [Ha].

The idea of using the construction of $\gr_m(A)$ to relate the study of a
general Noetherian local ring $(A, m)$ to the theory of ideals in a polynomial
ring over a field was one of the crucial ideas introduced by Krull in his
article `Dimension theory of local rings' [6], a work of monumental
significance for the theory of Noetherian rings. If $m$ is generated by $r$
elements then $\gr_m(A)$ is of the form $k[X_1, \ldots, X_r]/I$, where $k = A/m$ and $I$ is a
homogeneous ideal. However, the Hilbert function of this graded ring was
first used in the study of the multiplicity of $A$ by P. Samuel (1951).

\subsection*{Samuel functions}

In a little more generality, let $A$ be a Noetherian semilocal ring, and $m$
the Jacobson radical of $A$. If $I$ is an ideal of $A$ such that for some $\nu > 0$
we have $m^\nu \subset I \subset m$, we call $I$ an \emph{ideal of definition}\index{ideal of definition}; the $I$-adic and $m$-adic
topologies then coincide, so that `ideal of definition' means `ideal defining
the $m$-adic topology'. Let $M$ be a finite $A$-module. If we set
\[
  \gr_I(M) = \bigoplus_{n \ge 0} I^nM/I^{n+1}M
\]
then $\gr_I(M)$ is in a natural way a graded module over $\gr_I(A) = \bigoplus I^n/I^{n+1}$.
For brevity write $\gr_I(A) = A^*$ and $\gr_I(M) = M^*$. Then the ring $A^*_0 = A/I$ is
Artinian, and if $I = \sum_{i=1}^{r} x_iA$, and $\bar{x}_i$ is the image of $x_i$ in $I/I^2$, then
$A^* = A^*_0[\bar{x}_1, \ldots, \bar{x}_r]$. If also $M = \sum A\omega_i$ then $M^* = \sum A^*\bar{\omega}_i$ (where $\bar{\omega}_i$ is the
image of $\omega_i$ in $M_0^* = M/IM$), so that we can apply Theorem 2 and its
corollary to $M^*$. Noting that $l(M_n^*) = l(I^nM/I^{n+1}M)$ (where on the left-hand
side $l$ is the length as an $A^*_0$-module, on the right-hand side as an $A$-module),
we have
\[
  \sum_{n=0}^{\infty} l(M_n^*)t^n = \sum_{n=0}^{\infty} l(I^nM/I^{n+1}M)t^n.
\]
We now set $\chi_M^I(n) = l(M/I^{n+1}M)$. In particular we abbreviate $\chi_A^I(n)$ to
$\chi_A(n)$, and call it the \emph{Samuel function}\index{Samuel function} of the $A$-module $M$.

Repeatedly using the well-known formula
\[
  \binom{n}{r} = \binom{n-1}{r} + \binom{n-1}{r-1}
\]
we get
\[
  l(M/I^{n+1}M) = \sum_{k=0}^{n} l(I^kM/I^{k+1}M),
\]
so that from formula $(*)$ on p. 95 we get
\[
  \chi_M^I(n) = a_dn^d + a_{d-1}n^{d-1} + \cdots + a_0 \quad \text{with } a_i \in \Z.
\]
When $n \ge s$ this is a polynomial in $n$ of degree $d$. This degree $d$ is determined by $M$, and does not depend on $I$; to see this, if $I$ and $J$ are
both ideals of definition of $A$ then there exist natural numbers $a$ and $b$
such that $I^a \subset J$, $J^b \subset I$, so that
\[
  \chi_M^I(an + a - 1) \ge \chi_M^J(n) \quad \text{and} \quad \chi_M^J(bn + b - 1) \ge \chi_M^I(n).
\]
We thus write $d = d(M)$. It is natural to think of $d(M)$ as a measure of the size
of $M$.

\begin{theorem}\label{thm:13.3}
Let $A$ be a semilocal Noetherian ring, and $0 \to M' \to
M \to M'' \to 0$ an exact sequence of finite $A$-modules; then
\[
  d(M) = \max(d(M'), d(M'')).
\]
If $I$ is any ideal of definition of $A$, then $\chi_M^I$ and $\chi_{M''}^I$ have the same
leading coefficient.
\end{theorem}

\begin{proof}
We can assume $M'' = M/M'$. Then since $M''/I^nM'' = M/(M' + I^nM)$
we have
\[
  l(M/I^nM) = l(M/M' + I^nM) + l(M' + I^nM/I^nM)
\]
\[
  = l(M''/I^nM'') + l(M'/M' \cap I^nM).
\]
Thus setting $q(n) = l(M'/M' \cap I^{n+1}M)$, we have $\chi_M^I = \chi_{M''}^I + q$. Since
moreover both $\chi_{M''}^I$ and $q$ take on only positive values, $d(M)$ coincides with
whichever is the greater of $d(M'')$ and $\deg q$. However, by the Artin--Rees
lemma, there is a $c > 0$ such that
\[
  I^{n+1}M \cap M' \subset M' \cap I^{n+1}M \subset I^{n+1-c}M',
\]
and hence
\[
  \chi_{M''}^I(n) \ge q(n) \ge \chi_{M'}^I(n - c);
\]
therefore $q$ and $\chi_{M'}^I$ have the same leading coefficient.
\end{proof}

We now define a further measure $\delta(M)$ of the size of $M$: let $\delta(M)$ be the
smallest value of $n$ such that there exist $x_1, \ldots, x_n \in m$ for which $l(M/x_1M +
\cdots + x_nM) < \infty$. When $l(M) < \infty$, we interpret this as $\delta(M) = 0$. If $I$ is any
ideal of definition of $A$ then $l(M/IM) < \infty$, so that $\delta(M) \le$ number of
generators of $I$. Conversely, in the case that $A$ is a local ring and $M = A$,
then $l(A/I) < \infty$ implies that $I$ is an $m$-primary ideal. Therefore in this case
$\delta(A)$ is the minimum of the number of generators of $m$-primary ideals.

We have now arrived at the fundamental theorem of dimension theory.

\begin{theorem}\label{thm:13.4}
Let $A$ be a semilocal Noetherian ring and $M$ a finite
$A$-module; then we have
\[
  \dim M = d(M) = \delta(M).
\]
\end{theorem}

\begin{proof}
\textbf{Step 1.} Each of $d(M)$ and $\delta(M)$ are finite, but the finiteness of $\dim M$
has not yet been established. First of all, let us prove that $d(A) \ge \dim A$
for the case $M = A$, by induction on $d(A)$. Set $m = \rad(A)$. If $d(A) = 0$ then
$l(A/m^n)$ is constant for $n \gg 0$, so that for some $n$ we have $m^n = m^{n+1}$ and
by NAK, $m^n = 0$. Hence any prime ideal of $A$ is maximal and $\dim A = 0$.

Next suppose that $d(A) > 0$; if $\dim A = 0$ then we're done. If $\dim A > 0$,
consider a strictly increasing sequence $p_0 \subset p_1 \subset \cdots \subset p_e$ of prime ideals
of $A$, choose some element $x \in p_1 - p_0$, and set $B = A/(p_0 + xA)$; then by
the previous theorem applied to the exact sequence
\[
  0 \to A/p_0 \xrightarrow{x} A/p_0 \to B \to 0,
\]
we have $d(B) < d(A)$, and so by induction
\[
  \dim B \le d(B) \le d(A) - 1.
\]
(The values of $d(B)$ and $\dim B$ are independent of whether we consider $B$ as
an $A$-module or as a $B$-module, as is clear from the definitions.) In $B$, the
image of $p_1 \subset \cdots \subset p_e$ provides a chain of prime ideals of length $e - 1$, so that
\[
  e - 1 \le \dim B \le d(A) - 1;
\]
hence $e < d(A)$. Since this holds for any chain of prime ideals of $A$, this
proves $\dim A \le d(A)$. For general $M$, by Theorem~\ref{thm:6.4}, there are submodules
$M_i$ such that $0 = M_0 \subset M_1 \subset \cdots \subset M_r = M$ with
\[
  M_i/M_{i-1} \cong A/p_i, \quad \text{and} \quad p_i \in \Spec A.
\]
Since for an exact sequence $0 \to M' \to M \to M'' \to 0$ of finite $A$-
modules we have
\[
  \Supp(M) = \Supp(M') \cup \Supp(M'')
\]
and
\[
  \dim M = \max(\dim M', \dim M''),
\]
it is easy to see that
\[
  d(M) = \max\{d(A/p_i)\} \ge \max\{\dim(A/p_i)\} = \dim M.
\]

\textbf{Step 2.} We show that $\delta(M) \ge d(M)$. If $\delta(M) = 0$ then $l(M) < \infty$ so that $\chi_M(n)$
is bounded, hence $d(M) = 0$. Next suppose that $\delta(M) = s > 0$, choose
$x_1, \ldots, x_s \in m$ such that $l(M/x_1M + \cdots + x_sM) < \infty$, and set $M_i = M/x_1M +
\cdots + x_iM$; then clearly $\delta(M_i) = \delta(M) - i$. On the other hand,
\[
  l(M_i/m^\nu M_i) = l(M/x_iM + m^\nu M)
\]
\[
  = l(M/m^\nu M) - l(x_iM/x_iM \cap m^\nu M)
\]
\[
  = l(M/m^\nu M) - l(M/(m^\nu M:x_i))
\]
\[
  \ge l(M/m^\nu M) - l(M/m^{\nu-1}M),
\]
so that $d(M_i) \ge d(M) - 1$. Repeating this, we get $d(M_s) \ge d(M) - s$, but since
$\delta(M_s) = 0$ we have $d(M_s) = 0$, so that $s \ge d(M)$.

\textbf{Step 3.} We show that $\dim M \ge \delta(M)$, by induction on $\dim M$. If $\dim M = 0$
then $\Supp(M) \subset \text{m-Spec } A = V(m)$ so that for large enough $n$ we have
$m^n \subset \ann M$, and $l(M) < \infty$, therefore $\delta(M) = 0$. Next suppose that
$\dim M > 0$, and let $p_i$ for $1 \le i \le t$ be the minimal prime divisors of $\ann(M)$
with $\coht p_i = \dim M$; then the $p_i$ are not maximal ideals, so do not contain
$m$. Hence we can choose $x \in m$ not contained in any $p_i$. Setting
$M_1 = M/xM$ we get $\dim M_1 < \dim M$. Therefore by the inductive hypo-
thesis $\delta(M_1) \le \dim M_1$; but obviously $\delta(M) \le \delta(M_1) + 1$, so that $\delta(M) \le
\dim M_1 + 1 \le \dim M$.
\end{proof}

\begin{theorem}\label{thm:13.5}
Let $A$ be a Noetherian ring, and $I = (a_1, \ldots, a_r)$ an ideal
generated by $r$ elements; then if $p$ is a minimal prime divisor of $I$ we have
$\htop p \le r$. Hence the height of a proper ideal of $A$ is always finite.
\end{theorem}

\begin{proof}
The ideal $IA_p \subset A_p$ is a primary ideal belonging to the maximal
ideal, so that $\htop p = \dim A_p = \delta(A_p) \le r$.
\end{proof}

\begin{remark}
Krull proved this theorem by induction on $r$; the case $r = 1$ is
then the hardest part of the proof. Krull called the $r = 1$ case the
\emph{principal ideal theorem} (Hauptidealsatz), and the whole of Theorem 5 is sometimes
known by this name. Here Theorem 5 is merely a corollary of Theorem 4,
but one can also deduce the statement $\dim M = \delta(M)$ of Theorem 4 from
it. As far as proving Theorem 5 is concerned, Krull's proof, which does
not use the Samuel function, is easier. For this proof, see [N1] or [K].
More elementary proofs of the principal ideal theorem can be found in
Rees [3] and Caruth [1].

The definition of height is abstract, and even when one can find a lower
bound, one cannot expect an upper bound just from the definition, so
that this theorem is extremely important. The principal ideal theorem
corresponds to the familiar and obvious-looking proposition of geo-
metrical and physical intuition (which is strictly speaking not always true)
that `adding one equation can decrease the dimension of the space of
solutions by at most one'.
\end{remark}

\begin{theorem}\label{thm:13.6}
Let $P$ be a prime ideal of height $r$ in a Noetherian ring $A$.
Then
\begin{enumerate}[(i)]
  \item $P$ is a minimal prime divisor of some ideal $(a_1, \ldots, a_r)$ generated by
        $r$ elements;
  \item if $b_1, \ldots, b_s \in P$ we have $\htop P/(b_1, \ldots, b_s) \ge r - s$;
  \item if $a_1, \ldots, a_r$ are as in (i) we have
        \[
          \htop P/(a_1, \ldots, a_i) = r - i \quad \text{for } 1 \le i \le r.
        \]
\end{enumerate}
\end{theorem}

\begin{proof}
(i) $A_P$ is an $r$-dimensional local ring, so that by Theorem 4 we can
choose $r$ elements $a_1, \ldots, a_r \in PA_P$ such that $(a_1, \ldots, a_r)A_P$ is $PA_P$-primary.
Each $a_i$ is of the form an element of $P$ times a unit of $A_P$, so that without
loss of generality we can assume that $a_i \in P$. Then $P$ is a minimal prime
divisor of $(a_1, \ldots, a_r)A$.

(ii) Set $\bar{A} = A/(b_1, \ldots, b_s)$, $\bar{P} = P/(b_1, \ldots, b_s)$ and $\htop P = t$. Then by (i),
there exist $c_1, \ldots, c_t \in P$ such that $P$ is a minimal prime divisor of
$(c_1, \ldots, c_t)A$. Then $\bar{P}$ is a minimal prime divisor of $(\bar{b}_1, \ldots, \bar{b}_s, \bar{c}_1, \ldots, \bar{c}_t)$,
and hence $r \le s + t$ by Theorem 5.

(iii) In $\bar{A} = A/(a_1, \ldots, a_i)$, the ideal $\bar{P} = P/(a_1, \ldots, a_i)$ is a minimal
prime divisor of $(\bar{a}_{i+1}, \ldots, \bar{a}_r)$, hence
\[
  \htop P/(a_1, \ldots, a_i) \le r - i.
\]
The opposite inequality was proved in (ii).
\end{proof}

\begin{theorem}\label{thm:13.7}
Let $A = \bigoplus_{n \ge 0} A_n$ be a Noetherian graded ring.
\begin{enumerate}[(i)]
  \item If $I$ is a homogeneous ideal and $P$ is a prime divisor of $I$ then $P$ is also
        homogeneous.
  \item If $P$ is a homogeneous prime ideal of height $r$ then there exists a
        sequence
        \[
          P = P_r \supset P_{r-1} \supset \cdots \supset P_0
        \]
        of length $r$ consisting of homogeneous
        prime ideals.
\end{enumerate}
\end{theorem}

\begin{proof}
(i) $P$ can be expressed in the form $P = \ann(x)$ for a suitable element
$x$ of the graded $A$-module $A/I$. Let $a \in P$, and let $x = x_0 + x_1 + \cdots + x_n$
and $a = a_p + a_{p+1} + \cdots + a_q$ be decompositions into homogeneous terms.
Then since $ax = 0$,
\[
  a_px_0 = 0, \quad a_px_1 + a_{p+1}x_0 = 0, \quad a_px_2 + a_{p+1}x_1 + a_{p+2}x_0 = 0, \ldots,
\]
from which we get $a_p^2x_0 = 0$, $a_p^2x_1 = 0, \ldots$, and finally $a_p^{n+1}x = 0$. It follows
that $a_p^{n+1} \in P$, but since $P$ is prime, $a_p \in P$. Thus $a_{p+1} + \cdots + a_q \in P$, so
that in turn $a_{p+1} \in P$. Proceeding in the same way, we see that all the homo-
geneous terms of $a$ are in $P$, so that $P$ is a homogeneous ideal.

(ii) First of all note that we can assume that $A$ is an integral domain.
To see this, if we take a chain $P = P_r \supset \cdots \supset P_0$ of prime ideals of length
$r$ then $P_0$ is a minimal prime divisor of $(0)$, and so by (i) is a homogeneous
ideal; so we can replace $A$ by $A/P_0$. Now choose a homogeneous element
$0 \neq b_r \in P$; then by Theorem 6, $\htop (P/b_rA) = r - 1$, and so there is a minimal
prime divisor $Q$ of $b_rA$ such that $\htop (P/Q) = r - 1$; since $Q \neq (0)$, it is a
height $1$ homogeneous prime ideal. By the inductive hypothesis on $r$
applied to $P/Q$ there exists a chain $P = P_r \supset P_{r-1} \supset \cdots \supset P_1 = Q$ of homo-
geneous prime ideals of length $r - 1$, and adding on $P_0 = (0)$ we get a chain
of length $r$.
\end{proof}

Let us investigate more closely the relation between local rings and
graded rings.

\begin{theorem}\label{thm:13.8}
Let $k$ be a field, and $R = k[\xi_1, \ldots, \xi_r]$ a graded ring
generated by elements $\xi_1, \ldots, \xi_r$ of degree $1$; set $M = \sum \xi_iR$, $A = R_M$, and
$m = MA_M$.
\begin{enumerate}[(i)]
  \item Let $\chi$ be the Samuel function of the local ring $A$, and $\varphi$ the Hilbert
        function of the graded ring $R$; then $\varphi(n) = \chi(n) - \chi(n - 1)$;
  \item $\dim R = \htop M = \dim A = \deg \varphi + 1$;
  \item $\gr_m(A) \cong R$ as graded rings.
\end{enumerate}
\end{theorem}

\begin{proof}
$M$ is a maximal ideal of $R$ so that
\[
  m^n/m^{n+1} \cong M^n/M^{n+1} \cong R_n;
\]
hence $\chi(n) - \chi(n - 1) = l(m^n/m^{n+1}) = l(R_n) = \varphi(n)$, and so $\dim A = \deg \chi = 1 +
\deg \varphi$. Then since $A = R_M$, we have $\dim A = \htop M$. After this it is enough
to prove that $\dim R = \htop M$. First of all, assume that $R$ is an integral
domain, so that by Example 3 in the section on Hilbert functions and by
Theorem~\ref{thm:5.6}, we have
\[
  1 + \deg \varphi \ge \trdeg_k R = \dim R \ge \htop M;
\]
putting this together with $\htop M = \dim A = 1 + \deg \varphi$, we get $\dim R = \htop M$.
Next for general $R$, let $P_1, \ldots, P_s$ be the minimal prime ideals of $R$; then by
Theorem 7, these are all homogeneous ideals, and each $R/P_i$ is a graded
ring. Choosing $P_1$ such that $\dim R = \dim R/P_1$ and using the above result,
we get
\[
  \dim R = \dim R/P_1 = \htop M/P_1 \le \htop M \le \dim R,
\]
so that $\dim R = \htop M$ as required. We have $R_n \subset M^n \subset m^n$ with $m^n/m^{n+1} \cong
R_n$, and so taking an element $x$ of $R_n$ into its image in $m^n/m^{n+1}$ we obtain
a canonical one-to-one map $R \cong \gr_m A$, and it is clear from the definition
that this is a ring isomorphism.
\end{proof}

\begin{theorem}\label{thm:13.9}
Let $(A, m, k)$ be a Noetherian local ring, and set $G = \gr_m A$;
then $\dim A = \dim G$.
\end{theorem}

\begin{proof}
Letting $\varphi$ be the Hilbert polynomial of $G$, we have $\dim A =
1 + \deg \varphi$ (by Theorem 4), and by the previous theorem this is equal to
$\dim G$.
\end{proof}

In fact the following more general theorem holds: for $I$ a proper ideal in a
Noetherian local ring $A$, set $G = \gr_I(A)$; then $\dim A = \dim G$. This will be
proved a little later (Theorem~\ref{thm:15.7}).

% ============ 习题 ============
\section*{Exercises to \S\ref{sec:13}}
\emph{Prove the following propositions.}

\begin{exercise}\label{ex:13.1-ex}
Let $R = R_0 + R_1 + \cdots$ be a graded ring, and $u$ a unit of $R_0$. Then the map $T_u$,
defined by $T_u(x_0 + x_1 + \cdots + x_n) = x_0 + x_1u + \cdots + x_nu^n$ (where $x_i \in R_i$) is
an automorphism of $R$. If $R_0$ contains an infinite field $k$, then an ideal $I$ of $R$
is homogeneous if and only if $T_u(I) = I$ for every $u \in k$.
\end{exercise}

\begin{exercise}\label{ex:13.2-ex}
Let $R = R_0 + R_1 + \cdots$ be a graded ring, $I$ an ideal of $R$ and $t$ an
indeterminate over $R$. Set $R' = R[t, t^{-1}]$ and consider $R'$ as a graded ring
where $t$ has degree $0$ (that is, $R'_n = R_n[t, t^{-1}]$). Then an ideal $I$ of $R$ is
homogeneous if and only if $T_u(IR') = IR'$.
\end{exercise}

\begin{exercise}\label{ex:13.3-ex}
Let $A$ be a Noetherian ring having an embedded associated prime. If $a \in A$
is a non-zero divisor satisfying $\bigcap_{n \ge 1} a^nA = (0)$, then $A/(a)$ also has an
embedded associated prime.
\end{exercise}

\begin{exercise}\label{ex:13.4-ex}
Let $R = \bigoplus_{n \in \Z} R_n$ be a $\Z$-graded ring. For an ideal $I$ of $R$, let $I^*$ denote the
greatest homogeneous ideal of $R$ contained in $I$, that is the ideal of $R$
generated by all the homogeneous elements of $I$.
\begin{enumerate}[(i)]
  \item If $P$ is prime so is $P^*$.
  \item If $P$ is a homogeneous prime ideal and $Q$ is a $P$-primary ideal then $Q^*$
        is again $P$-primary.
\end{enumerate}
\end{exercise}

\begin{exercise}\label{ex:13.5-ex}
Let $R$ be a $\Z$-graded integral domain; write $S$ for the multiplicative set
consisting of all non-zero homogeneous elements of $R$. Then $R_S$ is a graded
ring, and its component of degree $0$ is a field $(R_S)_0 = K$; if $R \neq R_S$ then $R_S
\cong K[X, X^{-1}]$, where the degree of $X$ is the greatest common divisor of the
degrees of elements of $S$.
\end{exercise}

\begin{exercise}\label{ex:13.6-ex}
Let $R$ be an $\N$-graded ring and $P$ an inhomogeneous prime ideal of $R$; then
there are no prime ideals contained between $P^*$ and $P$. If $\htop P < \infty$ then
$\htop P = \htop P^* + 1$ (Matijevic--Roberts [1]).
\end{exercise}
